%------------------------------------------------------------------------------%
% Luis A. D'Afonseca
%
%------------------------------------------------------------------------------%
\documentclass[a4paper,12pt,fleqn]{article}

\usepackage{style_avaliacao}
\usepackage{systeme}

\ShowAnswers

\newcommand{\D}[2]{\frac{\partial {#1}}{\partial {#2}}}
\newcommand{\DS}[2]{\frac{\partial^2 {#1}}{\partial {#2}^2}}
\newcommand{\DM}[3]{\frac{\partial^2 {#1}}{\partial {#2} \partial {#3}}}

%------------------------------------------------------------------------------%
\begin{document}

\makeHeader{IS}{Prova Substitutiva -- Turma 8}{13/09/2024}

\textbf{O exercício correspondente a prova que será substituida vale 50 pontos}

% 01
%------------------------------------------------------------------------------%
\exercise{25}
Calcule o volume do sólido formado pela rotação da região entre as curvas
\[
y = x+2 \qquad
y = x^2
\]
em torno da reta \(x=-1\).

\begin{answer}
  Vamos usar cascas cilíndricas, para isso precisamos encontrar o raio e a
  altura da casca cilindica
  \[
    r(x) = x - (-1) = x + 1
  \]
  \[
    h(x) = (x + 2) - x^2 = 2 + x - x^2
  \]
  O intervalo de integração, \([a, b]\), vai ser entre os pontos
  comuns entre as curvas \(y=x+2\) e \(y=x^2\),
  ou seja \(a=0\) e
  \begin{align*}
    (x + 2) & = x^2 \\
    x^2 - x - 2 = 0
  \end{align*}
  \[
    \Delta
    = b^2-4ac
    = (-1)^2 - 4(1)(-2)
    = 9
    \qquad\qquad
    \sqrt{\Delta} = 3
  \]
  \[
    x
    = \frac{-b\pm\sqrt{\Delta}}{2a}
    = \frac{-(-1)\pm 3}{2\times 1}
    = \frac{1\pm 3}{2}
  \]
  portanto
  \[
    x_1 = \frac{1-3}{2} = \frac{-2}{2} = -1
    \qquad
    x_2 = \frac{1+3}{2} = \frac{4}{2} = 2
  \]
  O intervalo de integração é \([-1, 2]\)

  \vspace{\baselineskip}
  Calculando o volume
    \begin{align*}
      V
      & = \int_a^b 2\pi \;r(x) \;h(x)\; dx  \\[2mm]
      & = 2\pi \int_{-1}^2 (x + 1) (2 + x - x^2) dx  \\[2mm]
      & = 2\pi \int_{-1}^2 2x + x^2 - x^3 + 2 + x - x^2  dx \\[2mm]
      & = 2\pi \int_{-1}^2 2 + 3x -x^3 dx \\[2mm]
      & = 2\pi \left(2x + 3\frac{x^2}{2} - \frac{x^4}{4}\right) \evalat{-1}{2} \\[2mm]
      & = 2\pi \left(
                 \left[2\times 2 + 3\frac{2^2}{2} - \frac{2^4}{4}\right] -
                 \left[2(-1) + 3\frac{(-1)^2}{2} - \frac{(-1)^4}{4}\right]
               \right) \\[2mm]
      & = 2\pi \left(
                 \left[4 + 6 - 4\right] -
                 \left[-2 + \frac{3}{2} - \frac{1}{4}\right]
               \right) \\[2mm]
      & = 2\pi \left( 6 - \frac{-8 + 6 - 1}{4} \right) \\[2mm]
      & = 2\pi \left( \frac{24+3}{4} \right) \\[2mm]
      & = \frac{27\times 2\pi}{4} \\[2mm]
      & = \frac{27\pi}{2}
    \end{align*}
  \clearpage
\end{answer}

% 02
%------------------------------------------------------------------------------%
\exercise{25}
Calcule a integral
\(
  \int_1^e x^3\ln(x) dx
\)

\begin{answer}
  Para calcular a primitiva
  \[
    F(x) = \int x^3\ln(x) dx
  \]
  vamos usar integração por partes com
  \[
    u = \ln(x) \qquad du = \frac{dx}{x} \qquad\qquad
    dv = x^3dx \qquad v = \frac{x^4}{4}
  \]
  obtendo
  \begin{align*}
    F(x)
    & = \int x^3\ln(x) dx \\[2mm]
    & = \ln(x)\frac{x^4}{4} - \int \frac{x^4}{4} \frac{dx}{x} \\[2mm]
    & = \frac{x^4\ln(x)}{4} - \frac{1}{4} \int x^3 dx  \\[2mm]
    & = \frac{x^4\ln(x)}{4} - \frac{1}{4} \frac{x^4}{4} + C  \\[2mm]
    & = \frac{x^4}{4}\left(\ln(x) - \frac{1}{4}\right) + C
  \end{align*}
  Agora podemos calcular
  \begin{align*}
    I
    & = \int_1^e x^3\ln(x) dx \\[2mm]
    & = F(x)\evalat{1}{e} \\[2mm]
    & = \left(\frac{x^4}{4}\left(\ln(x) - \frac{1}{4}\right)\right)\evalat{1}{e} \\[2mm]
    & = \left(\frac{e^4}{4}\left(\ln(e) - \frac{1}{4}\right)\right)
      - \left(\frac{1^4}{4}\left(\ln(1) - \frac{1}{4}\right)\right) \\[2mm]
    & = \frac{e^4}{4}\left(1 - \frac{1}{4}\right)
      - \frac{1}{4}\left(0 - \frac{1}{4}\right) \\[2mm]
    & = \frac{e^4}{4}\;\frac{3}{4}
      + \frac{1}{4}\frac{1}{4} \\[2mm]
    & = \frac{3e^4+ 1}{16}
  \end{align*}
  \clearpage
\end{answer}

% 03
%------------------------------------------------------------------------------%
\exercise{25}
Encontre a série de Taylor, centrada em zero, da função \(e^{-x}\)

\begin{answer}
\end{answer}

%------------------------------------------------------------------------------%
\end{document}
%------------------------------------------------------------------------------%
